% Fate's Edge SRD --- Character Creation (Section 03)

\section{Character Creation}

\subsection{Starting Build Points}

Characters begin with \textbf{32 XP} to allocate during initial creation. This represents a balanced baseline for competent starting characters.

\subsection{Enhanced Starting Builds}

Players may exceed the standard 32 XP build through narrative engagement:

\begin{itemize}
  \item \textbf{Bonds:} Up to two player-defined mutual bonds may be taken for +2 XP total. (See Section~\texttt{[MISSING REF: sec:coordination]} for bond-driven Boon rules.)
  \item \textbf{Complications:} Up to two initial complications may be accepted for +4 XP total. Each complication causes scenes to start with +1 banked SB per character with that complication until the complication has been resolved in play.
\end{itemize}

This allows a maximum starting build of \textbf{36 XP}. Players are encouraged to aim for 32 XP and use bonds and complications to enhance characterization rather than pure mechanical optimization.

\subsection{Step-by-Step Creation}

\subsubsection{Step 1: Concept and Culture}
Describe your character in a sentence. Choose a cultural background (e.g., Aeler, Ykrul, Vilikari, Lethai). Cultural affinities provide narrative edges described in the full cultural guides.

\subsubsection{Step 2: Attributes}
Distribute 1, 2, 2, 3 among the four attributes (minimum 1, maximum 3 at start). Cost to raise attributes later is new rating × 3 XP.

\begin{tabular}{@{}ll@{}}
\toprule
\textbf{Attribute} & \textbf{Description} \\
\midrule
Body & Strength, endurance, physical resilience. \\
Wits & Perception, speed, quick thinking. \\
Spirit & Willpower, focus, connection to the intangible. \\
Presence & Charisma, command, force of personality. \\
\bottomrule
\end{tabular}

\subsubsection{Step 3: Skills}
Distribute 6 points among skills (maximum 2 in any skill at start). Cost to raise skills later is new level × 2 XP.

Core skills include:
\begin{multicols}{2}
\begin{itemize}
  \item Melee
  \item Ranged
  \item Brawl
  \item Athletics
  \item Stealth
  \item Endurance
  \item Craft
  \item Survival
  \item Sway
  \item Command
  \item Deception
  \item Performance
  \item Insight
  \item Lore
  \item Investigation
  \item Medicine
  \item Arcana
  \item Ritual
\end{itemize}
\end{multicols}

\subsubsection{Step 4: Talents}
Spend XP on talents. Recommended starting talents for common archetypes:

\begin{tabular}{@{}ll@{}}
\toprule
\textbf{Archetype} & \textbf{Recommended Talents} \\
\midrule
Warrior & Weapon Master (5 XP), Effortless Charge (4 XP), Conditioning (4 XP) \\
Rogue & Keen Senses (2 XP), Trackless Step (2 XP), Backstab (6 XP) \\
Caster (Freeform) & Spellcraft (6 XP) \\
Runekeeper & Familiar (2 XP) + Codex (4 XP) \\
Invoker & Patron's Symbol (4 XP) \\
Summoner & Pact-Whisperer (2 XP) + Lesser Pactwright (2 XP) \\
Support & Field Dressing (2 XP), Medic's Hand (3 XP), Inspire (3 XP) \\
\bottomrule
\end{tabular}

\subsubsection{Step 5: Equipment}
Choose starting equipment based on concept. See Section~\texttt{[MISSING REF: sec:equipment]} for weapon, armor, and gear tables.

\subsubsection{Step 6: Bonds}
Define 1--2 mutual bonds with other player characters. Each bond gives:

\begin{itemize}
  \item Once per session, when you act meaningfully on the bond and provide an Intricate Description, gain 1 Boon.
  \item Bonded allies may gift 2 Boons per scene to each other (instead of the normal 1).
\end{itemize}

A bond should include:
\begin{itemize}
  \item A shared history or obligation.
  \item What each character gains from the bond.
  \item A potential source of tension.
\end{itemize}

\subsubsection{Step 7: Complications (Optional)}
Choose 0--2 complications. Each gives +2 XP at start. Complications should be narrative hooks, not purely mechanical penalties. Examples:

\begin{itemize}
  \item \textbf{Debt to a Patron:} You owe Obligation segments to a Patron before play begins.
  \item \textbf{Hunted:} A faction or individual seeks you.
  \item \textbf{Forbidden Knowledge:} You carry a secret that could destroy you if revealed.
  \item \textbf{Cursed Lineage:} Ancestral curse manifests at dramatic moments.
  \item \textbf{Oath-Bound:} A geas compels certain actions; breaking it has steep costs.
\end{itemize}

For each complication taken, at the start of each scene the GM gains +1 SB until the complication is resolved in play.

\subsubsection{Step 8: Finalize}
Record starting Fatigue (0), Harm (0), Boons (0). Note any starting Assets or Followers (see Section~\texttt{[MISSING REF: sec:followers-assets]}).

\subsection{Quick Build Templates}

\subsubsection{Fighter (32 XP)}
\begin{tabular}{@{}ll@{}}
Body 3, Wits 2, Spirit 2, Presence 2 \\
Skills: Melee 2, Athletics 2, Endurance 1, Command 1 \\
Talents: Weapon Master (5 XP), Conditioning (4 XP), Effortless Charge (4 XP) \\
Equipment: Medium weapon, medium armor \\
Bond: One mutual bond \\
Complication: None \\
\end{tabular}

\subsubsection{Rogue (32 XP)}
\begin{tabular}{@{}ll@{}}
Body 2, Wits 4, Spirit 2, Presence 2 \\
Skills: Stealth 2, Subterfuge 2, Sway 2, Tinker 1 \\
Talents: Keen Senses (2 XP), Silk Palms (2 XP), Backstab (6 XP) \\
Equipment: Light weapons, light armor, lockpicks \\
Bond: One mutual bond \\
Complication: None \\
\end{tabular}

\subsubsection{Caster (32 XP)}
\begin{tabular}{@{}ll@{}}
Body 2, Wits 4, Spirit 3, Presence 1 \\
Skills: Arcana 2, Lore 2, Insight 1, Investigation 1 \\
Talents: Spellcraft (6 XP) \\
Equipment: Staff, robes, focus (crystal, runestone, or inkpot) \\
Bond: One mutual bond \\
Complication: None \\
\end{tabular}

\subsection{Learning Through Examples}

See how core mechanics apply to the QuickBuild Templates (Sections~\texttt{[MISSING REF: sec:quickbuild-fighter]}--\texttt{[MISSING REF: sec:quickbuild-caster]}):

\paragraph{Lockpick Under Pressure (Lyra the Rogue, Section~\texttt{[MISSING REF: sec:quickbuild-rogue]}):}
\begin{itemize}
  \item GM: ``Controlled, DV 3. On a miss, the guard hears you.''
  \item Lyra rolls [9, 6, 4, 2, 1] (Wits 4 + Tinker 2 = 6 dice) → 2 successes (DV 3 Partial) → \textbf{Gains 1 Boon}.
  \item The 1 gives GM 1 SB: ``The lock clicks open, but the last pin squeals. I spend the SB: the guard stops and turns his head. You have one beat to hide.''
\end{itemize}
→ \textit{Teaches:} Partial success = progress + Boon; 1s = GM currency for complications.

\paragraph{Arcane Backlash (Sera the Caster, Section~\texttt{[MISSING REF: sec:quickbuild-caster]}):}
\begin{itemize}
  \item Sera casts [FIRE][STRIKE] (DV 3) vs. a fire elemental → rolls [10, 1, 1] (Wits 4 + Arcana 2 = 6 dice) → 3 successes (10=2, 1s=2 SB) → \textbf{Success with SB} (GM spends 1 SB for complication).
  \item Backlash: Dominant element is Fire → GM describes: ``The bolt hits, but the flash blinds you---your next action is Desperate. Smoke triggers a hanging lantern; the room starts catching fire (Clock: Fire Spread [4]).''
\end{itemize}
→ \textit{Teaches:} 1s generate SB \emph{even on success}; elemental backlash coloring; SB spent on visible consequences.

\paragraph{Tip for New GMs:}
Run these examples at your table. Ask: ``What would you do with Lyra's Boon?'' or ``How does Sera's backlash change the scene?''

\subsection{Advanced Start: Tier II (80 XP)}

For campaigns beginning at higher power, characters may start with 80 XP. In addition to the 32 XP baseline, add:

\begin{itemize}
  \item +2 attribute points (raise any attributes; maximum 4).
  \item +8 skill points (maximum skill 3 at start; raise to 4 later).
  \item Additional talents worth approximately 48 XP.
  \item One Minor Asset (4 XP) or one Follower (Cap 2).
  \item One additional bond or complication (narrative only; no extra XP).
\end{itemize}

\subsection{Session Zero Tips}

\begin{itemize}
  \item Establish the table's tone and stakes; tie character goals to campaign fronts.
  \item Map Bonds among PCs; mark possible bond-driven Boon triggers.
  \item Seed 1--2 personal Complications per PC for early spotlight.
  \item Discuss Patron choices (if any) and what their omens look like.
  \item Review Safety Tools (Lines, Veils, X-Card) before play begins.
\end{itemize}

\subsection{Build Point Summary}

\begin{tabular}{@{}lcc@{}}
\toprule
\textbf{Element} & \textbf{Cost (XP)} & \textbf{Starting Max} \\
\midrule
Attribute (per point) & 3 × new rating & 3 \\
Skill (per point) & 2 × new level & 2 \\
Talent & varies & --- \\
Bond (mutual) & +1 XP each (2 total) & 2 \\
Complication & +2 XP each (2 total) & 2 \\
Minor Asset & 4 XP & 1 \\
Follower (Cap 2) & 2 XP & 1 \\
\bottomrule
\end{tabular}


\subsubsection{Step 4: Talents}
Spend XP on talents. Recommended starting talents for common archetypes:

\begin{tabular}{@{}ll@{}}
\toprule
\textbf{Archetype} & \textbf{Recommended Talents} \\
\midrule
Warrior & Weapon Master (5 XP), Effortless Charge (4 XP), Conditioning (4 XP) \\
Rogue & Keen Senses (2 XP), Trackless Step (2 XP), Backstab (6 XP) \\
Caster (Freeform) & Spellcraft (6 XP) \\
Runekeeper & Familiar (2 XP) + Codex (4 XP) \\
Invoker & Patron's Symbol (4 XP) \\
Summoner & Pact-Whisperer (2 XP) + Lesser Pactwright (2 XP) \\
Support & Field Dressing (2 XP), Medic's Hand (3 XP), Inspire (3 XP) \\
\bottomrule
\end{tabular}

\subsubsection{Step 5: Equipment}
Choose starting equipment based on concept. See Section~\texttt{[MISSING REF: sec:equipment]} for weapon, armor, and gear tables.

\subsubsection{Step 6: Bonds}
Define 1--2 mutual bonds with other player characters. Each bond gives:

\begin{itemize}
  \item Once per session, when you act meaningfully on the bond and provide an Intricate Description, gain 1 Boon.
  \item Bonded allies may gift 2 Boons per scene to each other (instead of the normal 1).
\end{itemize}

A bond should include:
\begin{itemize}
  \item A shared history or obligation.
  \item What each character gains from the bond.
  \item A potential source of tension.
\end{itemize}

\subsubsection{Step 7: Complications (Optional)}
Choose 0--2 complications. Each gives +2 XP at start. Complications should be narrative hooks, not purely mechanical penalties. Examples:

\begin{itemize}
  \item \textbf{Debt to a Patron:} You owe Obligation segments to a Patron before play begins.
  \item \textbf{Hunted:} A faction or individual seeks you.
  \item \textbf{Forbidden Knowledge:} You carry a secret that could destroy you if revealed.
  \item \textbf{Cursed Lineage:} Ancestral curse manifests at dramatic moments.
  \item \textbf{Oath-Bound:} A geas compels certain actions; breaking it has steep costs.
\end{itemize}

For each complication taken, at the start of each scene the GM gains +1 SB until the complication is resolved in play.

\subsubsection{Step 8: Finalize}
Record starting Fatigue (0), Harm (0), Boons (0). Note any starting Assets or Followers (see Section~\texttt{[MISSING REF: sec:followers-assets]}).

\subsection{Quick Build Templates}

\subsubsection{Fighter (32 XP)}
\begin{tabular}{@{}ll@{}}
Body 3, Wits 2, Spirit 2, Presence 2 \\
Skills: Melee 2, Athletics 2, Endurance 1, Command 1 \\
Talents: Weapon Master (5 XP), Conditioning (4 XP), Effortless Charge (4 XP) \\
Equipment: Medium weapon, medium armor \\
Bond: One mutual bond \\
Complication: None \\
\end{tabular}

\subsubsection{Rogue (32 XP)}
\begin{tabular}{@{}ll@{}}
Body 2, Wits 4, Spirit 2, Presence 2 \\
Skills: Stealth 2, Subterfuge 2, Sway 2, Tinker 1 \\
Talents: Keen Senses (2 XP), Silk Palms (2 XP), Backstab (6 XP) \\
Equipment: Light weapons, light armor, lockpicks \\
Bond: One mutual bond \\
Complication: None \\
\end{tabular}

\subsubsection{Caster (32 XP)}
\begin{tabular}{@{}ll@{}}
Body 2, Wits 4, Spirit 3, Presence 1 \\
Skills: Arcana 2, Lore 2, Insight 1, Investigation 1 \\
Talents: Spellcraft (6 XP) \\
Equipment: Staff, robes, focus (crystal, runestone, or inkpot) \\
Bond: One mutual bond \\
Complication: None \\
\end{tabular}

\subsection{Advanced Start: Tier II (80 XP)}

For campaigns beginning at higher power, characters may start with 80 XP. In addition to the 32 XP baseline, add:

\begin{itemize}
  \item +2 attribute points (raise any attributes; maximum 4).
  \item +8 skill points (maximum skill 3 at start; raise to 4 later).
  \item Additional talents worth approximately 48 XP.
  \item One Minor Asset (4 XP) or one Follower (Cap 2).
  \item One additional bond or complication (narrative only; no extra XP).
\end{itemize}

\subsection{Session Zero Tips}

\begin{itemize}
  \item Establish the table's tone and stakes; tie character goals to campaign fronts.
  \item Map Bonds among PCs; mark possible bond-driven Boon triggers.
  \item Seed 1--2 personal Complications per PC for early spotlight.
  \item Discuss Patron choices (if any) and what their omens look like.
  \item Review Safety Tools (Lines, Veils, X-Card) before play begins.
\end{itemize}

\subsection{Build Point Summary}

\begin{tabular}{@{}lcc@{}}
\toprule
\textbf{Element} & \textbf{Cost (XP)} & \textbf{Starting Max} \\
\midrule
Attribute (per point) & 3 × new rating & 3 \\
Skill (per point) & 2 × new level & 2 \\
Talent & varies & --- \\
Bond (mutual) & +1 XP each (2 total) & 2 \\
Complication & +2 XP each (2 total) & 2 \\
Minor Asset & 4 XP & 1 \\
Follower (Cap 2) & 2 XP & 1 \\
\bottomrule
\end{tabular}

